% Λύση του 14ου προβλήματος της ενότητας «Άλυτα Προβλήματα».
\subsection*{Λύση Προβλήματος 14 --- Μέγιστη γωνία θέασης}

Θέτουμε $\alpha$ τη γωνία ανύψωσης προς το κάτω άκρο του πίνακα και
$\beta$ τη γωνία ανύψωσης προς το πάνω άκρο του. Τότε
\[
  \theta=\beta-\alpha.
\]
Επειδή το κάτω άκρο βρίσκεται $1\;\si{m}$ και το πάνω άκρο
$3\;\si{m}$ ψηλότερα από το επίπεδο των ματιών, έχουμε
\[
  \ef\alpha=\frac{1}{x}
  \qquad\text{και}\qquad
  \ef\beta=\frac{3}{x}.
\]
Άρα
\begin{align*}
  \ef\theta
  &=\ef(\beta-\alpha)\\
  &=\frac{\ef\beta-\ef\alpha}
          {1+\ef\beta\,\ef\alpha}\\
  &=\frac{\frac{3}{x}-\frac{1}{x}}
          {1+\frac{3}{x^2}}\\
  &=\frac{2x}{x^2+3}.
\end{align*}

Μελετούμε τη συνάρτηση
\[
  f(x)=\frac{2x}{x^2+3},
  \qquad x>0.
\]
Παραγωγίζουμε:
\[
  f'(x)
  =\frac{2(x^2+3)-4x^2}{(x^2+3)^2}
  =\frac{2(3-x^2)}{(x^2+3)^2}.
\]
Για $x>0$ ο παρονομαστής είναι θετικός. Επομένως,
\[
\begin{array}{c|ccc}
  x & (0,\sqrt3) & \sqrt3 & (\sqrt3,+\infty)\\ \hline
  f'(x) & + & 0 & -\\
  f(x) & \nearrow & \text{μέγιστο} & \searrow
\end{array}
\]
Άρα η $f$ γίνεται μέγιστη για
\[
  \boxed{x=\sqrt3\;\si{m}}.
\]

Επειδή $0<\theta<\dfrac{\pi}{2}$, η συνάρτηση $\ef\theta$ είναι
γνησίως αύξουσα. Συνεπώς, η μεγιστοποίηση της $f(x)=\ef\theta$ είναι
ισοδύναμη με τη μεγιστοποίηση της γωνίας $\theta$.

Για $x=\sqrt3$ βρίσκουμε
\[
  \ef\theta_{\max}
  =\frac{2\sqrt3}{3+3}
  =\frac{\sqrt3}{3}.
\]
Επομένως,
\[
  \boxed{\theta_{\max}=\frac{\pi}{6}=30^\circ}.
\]
